\documentclass[11pt, a4paper]{article}
\usepackage[utf8]{inputenc}
\usepackage[margin=1in]{geometry}
\usepackage{listings}
\usepackage{xcolor}
\usepackage{tabularx}

% Configure SQL code block styling
\lstset{
    basicstyle=\ttfamily\small,
    keywordstyle=\color{blue}\bfseries,
    stringstyle=\color{red!70!black},
    commentstyle=\color{green!60!black}\itshape,
    frame=single,
    backgroundcolor=\color{gray!5},
    breaklines=true,
    showstringspaces=false
}

\title{Assignment1 Lab Report}
\author{Student Name}
\date{\today}

\begin{document}

\maketitle

\section{Assignment 1}
\subsection*{Q1: Constraint violation attempts}
\begin{lstlisting}[language=SQL]
-- PRIMARY KEY violation
INSERT INTO
  depts (deptcode, deptname)
VALUES
  ('CSE', 'Duplicate Dept');

-- CHECK violation: bdate >= '1997-01-01'
INSERT INTO
  students (rollno, name, bdate, deptcode, hostel, parent_inc)
VALUES
  (
    99999999,
    'Test Student',
    '1998-05-10',
    'CSE',
    1,
    100000.0
  );

-- CHECK violation: hostel >= 10
INSERT INTO
  students (rollno, name, bdate, deptcode, hostel, parent_inc)
VALUES
  (
    99999998,
    'Test Student 2',
    '1995-01-01',
    'CSE',
    15,
    100000.0
  );

-- FOREIGN KEY violation: deptcode
INSERT INTO
  students (rollno, name, bdate, deptcode, hostel, parent_inc)
VALUES
  (
    99999997,
    'Test Student 3',
    '1995-01-01',
    'XYZ',
    1,
    100000.0
  );
\end{lstlisting}
\vspace{0.3cm}
\textcolor{red}{\textbf{Error:} duplicate key value violates unique constraint "depts\_pkey"
DETAIL:  Key (deptcode)=(CSE) already exists.}\vspace{0.2cm}\newline
\textcolor{red}{\textbf{Error:} new row for relation "students" violates check constraint "students\_bdate\_check"
DETAIL:  Failing row contains (99999999, Test Student                  , 1998-05-10, CSE, 1, 100000.0).}\vspace{0.2cm}\newline
\textcolor{red}{\textbf{Error:} new row for relation "students" violates check constraint "students\_hostel\_check"
DETAIL:  Failing row contains (99999998, Test Student 2                , 1995-01-01, CSE, 15, 100000.0).}\vspace{0.2cm}\newline
\textcolor{red}{\textbf{Error:} insert or update on table "students" violates foreign key constraint "students\_deptcode\_fkey"
DETAIL:  Key (deptcode)=(XYZ) is not present in table "depts".}\vspace{0.2cm}
\newpage
\subsection*{Q2: Delete deptcode='CSE' (cascades to students)}
\begin{lstlisting}[language=SQL]
SELECT
  rollno,
  name
FROM
  students
WHERE
  TRIM(deptcode) = 'CSE';

DELETE FROM crs_regd
WHERE
  crs_cd IN (
    SELECT
      crs_code
    FROM
      crs_offrd
    WHERE
      crs_fac_cd IN (
        SELECT
          fac_code
        FROM
          faculty
        WHERE
          TRIM(fac_dept) = 'CSE'
      )
  );

DELETE FROM crs_regd
WHERE
  crs_rollno IN (
    SELECT
      rollno
    FROM
      students
    WHERE
      TRIM(deptcode) = 'CSE'
  );

DELETE FROM crs_offrd
WHERE
  crs_fac_cd IN (
    SELECT
      fac_code
    FROM
      faculty
    WHERE
      TRIM(fac_dept) = 'CSE'
  );

DELETE FROM faculty
WHERE
  TRIM(fac_dept) = 'CSE';

DELETE FROM depts
WHERE
  TRIM(deptcode) = 'CSE';

SELECT
  rollno,
  name
FROM
  students
WHERE
  TRIM(deptcode) = 'CSE';
\end{lstlisting}
\vspace{0.3cm}
\textbf{Output:}\vspace{0.2cm}
\textit{CSE Students Before Deletion: }
\begin{center}
\begin{tabular}{|l|l|}
\hline
\textbf{rollno} & \textbf{name} \\ \hline
92005001 & Riddhi Agarwal                 \\ \hline
92005002 & Priya Verma                    \\ \hline
92005003 & Rahul Das                      \\ \hline
92005004 & Sneha Gupta                    \\ \hline
\end{tabular}
\end{center}
\textit{Statement executed successfully.}\vspace{0.2cm}\newline
\textit{Statement executed successfully.}\vspace{0.2cm}\newline
\textit{Statement executed successfully.}\vspace{0.2cm}\newline
\textit{Statement executed successfully.}\vspace{0.2cm}\newline
\textit{Statement executed successfully.}\vspace{0.2cm}\newline
\textit{Statement executed successfully. (0 rows returned)}\vspace{0.2cm} \newline
\newpage
\subsection*{Q3: Courses offered by faculty 'dbp' and 'nls'}
\begin{lstlisting}[language=SQL]
SELECT
  crs_code,
  crs_name
FROM
  crs_offrd
WHERE
  TRIM(crs_fac_cd) IN ('dbp', 'nls');
\end{lstlisting}
\vspace{0.3cm}
\textbf{Output:}\vspace{0.2cm}
\begin{center}
\begin{tabular}{|l|l|}
\hline
\textbf{crs\_code} & \textbf{crs\_name} \\ \hline
CS101 & DBMS                                \\ \hline
CS102 & OS                                  \\ \hline
CS103 & Software Engg.                      \\ \hline
CS104 & MIS                                 \\ \hline
\end{tabular}
\end{center}
\newpage
\subsection*{Q4: Courses with full details offered by 'dbp'}
\begin{lstlisting}[language=SQL]
SELECT
  co.crs_code,
  co.crs_name,
  co.crs_credits,
  f.fac_name,
  f.fac_dept
FROM
  crs_offrd co
  JOIN faculty f ON co.crs_fac_cd = f.fac_code
WHERE
  TRIM(co.crs_fac_cd) = 'dbp';
\end{lstlisting}
\vspace{0.3cm}
\textbf{Output:}\vspace{0.2cm}
\begin{center}
\begin{tabular}{|l|l|l|l|l|}
\hline
\textbf{crs\_code} & \textbf{crs\_name} & \textbf{crs\_credits} & \textbf{fac\_name} & \textbf{fac\_dept} \\ \hline
CS103 & Software Engg.                      & 6.0 & Dr. D.B. Phatak                & CSE \\ \hline
CS102 & OS                                  & 5.0 & Dr. D.B. Phatak                & CSE \\ \hline
CS101 & DBMS                                & 4.0 & Dr. D.B. Phatak                & CSE \\ \hline
\end{tabular}
\end{center}
\newpage
\subsection*{Q5: Courses with credits between 4.0 and 6.0}
\begin{lstlisting}[language=SQL]
SELECT
  crs_code,
  crs_name,
  crs_credits
FROM
  crs_offrd
WHERE
  crs_credits BETWEEN 4.0 AND 6.0;
\end{lstlisting}
\vspace{0.3cm}
\textbf{Output:}\vspace{0.2cm}
\begin{center}
\begin{tabular}{|l|l|l|}
\hline
\textbf{crs\_code} & \textbf{crs\_name} & \textbf{crs\_credits} \\ \hline
CS101 & DBMS                                & 4.0 \\ \hline
CS102 & OS                                  & 5.0 \\ \hline
CS103 & Software Engg.                      & 6.0 \\ \hline
CS104 & MIS                                 & 4.5 \\ \hline
EE101 & Circuit Theory                      & 5.5 \\ \hline
IT101 & Data Structures                     & 4.0 \\ \hline
\end{tabular}
\end{center}
\newpage
\subsection*{Q6: Courses with credits greater than 6.5}
\begin{lstlisting}[language=SQL]
SELECT
  crs_code,
  crs_name,
  crs_credits
FROM
  crs_offrd
WHERE
  crs_credits > 6.5;
\end{lstlisting}
\vspace{0.3cm}
\textbf{Output:}\vspace{0.2cm}
\begin{center}
\begin{tabular}{|l|l|l|}
\hline
\textbf{crs\_code} & \textbf{crs\_name} & \textbf{crs\_credits} \\ \hline
EE102 & Signal Processing                   & 7.0 \\ \hline
PH106 & Physics                             & 8.0 \\ \hline
ME101 & Thermodynamics                      & 7.5 \\ \hline
\end{tabular}
\end{center}
\newpage
\section{Assignment 2}
\subsection*{Q1: Count students in CSE dept}
\begin{lstlisting}[language=SQL]
SELECT
  COUNT(*)
FROM
  students
WHERE
  TRIM(deptcode) = 'CSE';
\end{lstlisting}
\vspace{0.3cm}
\textbf{Output:}\vspace{0.2cm}
\begin{center}
\begin{tabular}{|l|}
\hline
\textbf{count} \\ \hline
4 \\ \hline
\end{tabular}
\end{center}
\newpage
\subsection*{Q2: Min, max and average marks of each course}
\begin{lstlisting}[language=SQL]
SELECT
  crs_cd,
  MIN(marks) AS min_marks,
  MAX(marks) AS max_marks,
  AVG(marks) AS avg_marks
FROM
  crs_regd
GROUP BY
  crs_cd;
\end{lstlisting}
\vspace{0.3cm}
\textbf{Output:}\vspace{0.2cm}
\begin{center}
\begin{tabular}{|l|l|l|l|}
\hline
\textbf{crs\_cd} & \textbf{min\_marks} & \textbf{max\_marks} & \textbf{avg\_marks} \\ \hline
PH106 & 55.00 & 90.00 & 74.8750000000000000 \\ \hline
CS104 & 91.00 & 91.00 & 91.0000000000000000 \\ \hline
CH103 & 60.00 & 88.00 & 72.0000000000000000 \\ \hline
CS101 & 30.00 & 92.00 & 56.0500000000000000 \\ \hline
ME101 & 65.00 & 71.00 & 68.0000000000000000 \\ \hline
EE101 & 70.00 & 82.00 & 75.2500000000000000 \\ \hline
CS102 & 72.00 & 88.00 & 77.7500000000000000 \\ \hline
IT101 & 76.00 & 82.00 & 79.0000000000000000 \\ \hline
CS103 & 68.00 & 68.00 & 68.0000000000000000 \\ \hline
EE102 & 60.00 & 85.00 & 72.5000000000000000 \\ \hline
\end{tabular}
\end{center}
\newpage
\subsection*{Q3: Total credits of courses registered by a student}
\begin{lstlisting}[language=SQL]
SELECT
  r.crs_rollno,
  s.name,
  SUM(o.crs_credits) AS total_credits
FROM
  crs_regd r
  JOIN crs_offrd o ON r.crs_cd = o.crs_code
  JOIN students s on r.crs_rollno = s.rollno
GROUP BY
  r.crs_rollno,
  s.name;
\end{lstlisting}
\vspace{0.3cm}
\textbf{Output:}\vspace{0.2cm}
\begin{center}
\begin{tabular}{|l|l|l|}
\hline
\textbf{crs\_rollno} & \textbf{name} & \textbf{total\_credits} \\ \hline
92005004 & Sneha Gupta                    & 9.0 \\ \hline
92005012 & Vikram Singh                   & 7.0 \\ \hline
92005103 & Pooja Reddy                    & 15.0 \\ \hline
92005011 & Neha Patil                     & 9.5 \\ \hline
92005001 & Riddhi Agarwal                 & 54.5 \\ \hline
92005102 & Deepak Mehta                   & 16.5 \\ \hline
92005003 & Rahul Das                      & 4.0 \\ \hline
92005010 & Ravi Kumar                     & 9.5 \\ \hline
92005100 & Suresh Nair                    & 20.0 \\ \hline
92005002 & Priya Verma                    & 9.0 \\ \hline
92005104 & Arun Bhat                      & 8.0 \\ \hline
92005101 & Kavita Jain                    & 4.0 \\ \hline
92005106 & Kiran Desai                    & 7.5 \\ \hline
\end{tabular}
\end{center}
\newpage
\subsection*{Q4: Number of students in each hostel whose department is CSE}
\begin{lstlisting}[language=SQL]
SELECT
  hostel,
  COUNT(*) AS num_students
FROM
  students
WHERE
  TRIM(deptcode) = 'CSE'
GROUP BY
  hostel;
\end{lstlisting}
\vspace{0.3cm}
\textbf{Output:}\vspace{0.2cm}
\begin{center}
\begin{tabular}{|l|l|}
\hline
\textbf{hostel} & \textbf{num\_students} \\ \hline
1 & 2 \\ \hline
2 & 1 \\ \hline
3 & 1 \\ \hline
\end{tabular}
\end{center}
\newpage
\subsection*{Q5: Hostel, rollno, parent\_inc of student with max(parent\_inc) in each hostel}
\begin{lstlisting}[language=SQL]
SELECT
  hostel,
  rollno,
  parent_inc
FROM
  (
    SELECT
      hostel,
      rollno,
      parent_inc,
      RANK() OVER (
        PARTITION BY
          hostel
        ORDER BY
          parent_inc DESC
      ) as rank
    FROM
      students
    WHERE
      hostel IS NOT NULL
  ) ranked_students
WHERE
  rank = 1;
\end{lstlisting}
\vspace{0.3cm}
\textbf{Output:}\vspace{0.2cm}
\begin{center}
\begin{tabular}{|l|l|l|}
\hline
\textbf{hostel} & \textbf{rollno} & \textbf{parent\_inc} \\ \hline
1 & 92005001 & 500000.0 \\ \hline
2 & 92005002 & 750000.0 \\ \hline
3 & 92005004 & 600000.0 \\ \hline
4 & 92005011 & 900000.0 \\ \hline
5 & 92005012 & 350000.0 \\ \hline
6 & 92005100 & 800000.0 \\ \hline
7 & 92005101 & 650000.0 \\ \hline
8 & 92005103 & 700000.0 \\ \hline
9 & 92005105 & 550000.0 \\ \hline
\end{tabular}
\end{center}
\newpage
\subsection*{Q6: Name and parent\_inc of students with ncome greater than rollno 92005010}
\begin{lstlisting}[language=SQL]
SELECT
  name,
  parent_inc
FROM
  students
WHERE
  parent_inc > (
    SELECT
      parent_inc
    FROM
      students
    WHERE
      rollno = 92005010
  );
\end{lstlisting}
\vspace{0.3cm}
\textbf{Output:}\vspace{0.2cm}
\begin{center}
\begin{tabular}{|l|l|}
\hline
\textbf{name} & \textbf{parent\_inc} \\ \hline
Riddhi Agarwal                 & 500000.0 \\ \hline
Priya Verma                    & 750000.0 \\ \hline
Sneha Gupta                    & 600000.0 \\ \hline
Neha Patil                     & 900000.0 \\ \hline
Ananya Roy                     & 550000.0 \\ \hline
Suresh Nair                    & 800000.0 \\ \hline
Kavita Jain                    & 650000.0 \\ \hline
Pooja Reddy                    & 700000.0 \\ \hline
Meera Iyer                     & 550000.0 \\ \hline
Kiran Desai                    & 480000.0 \\ \hline
\end{tabular}
\end{center}
\newpage
\subsection*{Q7: Marks of students who scored more than rollno 92005102 in CH103 and PH106}
\begin{lstlisting}[language=SQL]
SELECT
  r.crs_rollno,
  r.crs_cd,
  r.marks
FROM
  crs_regd r
WHERE
  TRIM(r.crs_cd) IN ('CH103', 'PH106')
  AND r.marks > (
    SELECT
      r2.marks
    FROM
      crs_regd r2
    WHERE
      r2.crs_rollno = 92005102
      AND r2.crs_cd = r.crs_cd
  );
\end{lstlisting}
\vspace{0.3cm}
\textbf{Output:}\vspace{0.2cm}
\begin{center}
\begin{tabular}{|l|l|l|}
\hline
\textbf{crs\_rollno} & \textbf{crs\_cd} & \textbf{marks} \\ \hline
92005001 & CH103 & 88.00 \\ \hline
92005001 & PH106 & 74.50 \\ \hline
92005012 & CH103 & 65.00 \\ \hline
92005100 & CH103 & 75.00 \\ \hline
92005100 & PH106 & 80.00 \\ \hline
92005103 & PH106 & 90.00 \\ \hline
\end{tabular}
\end{center}
\newpage
\section{Assignment 3}
\subsection*{Q1: Students (rollno, name, deptcode) registered for course EE101}
\begin{lstlisting}[language=SQL]
SELECT
  s.rollno,
  s.name,
  s.deptcode
FROM
  students s
  JOIN crs_regd cr ON s.rollno = cr.crs_rollno
WHERE
  TRIM(cr.crs_cd) = 'EE101';
\end{lstlisting}
\vspace{0.3cm}
\textbf{Output:}\vspace{0.2cm}
\begin{center}
\begin{tabular}{|l|l|l|}
\hline
\textbf{rollno} & \textbf{name} & \textbf{deptcode} \\ \hline
92005001 & Riddhi Agarwal                 & CSE \\ \hline
92005010 & Ravi Kumar                     & ELE \\ \hline
92005011 & Neha Patil                     & ELE \\ \hline
92005102 & Deepak Mehta                   & ETC \\ \hline
\end{tabular}
\end{center}
\newpage
\subsection*{Q2: Students (rollno, name) in ELE dept registered for course EE101}
\begin{lstlisting}[language=SQL]
SELECT
  s.rollno,
  s.name
FROM
  students s
  JOIN crs_regd cr ON s.rollno = cr.crs_rollno
WHERE
  TRIM(s.deptcode) = 'ELE'
  AND TRIM(cr.crs_cd) = 'EE101';
\end{lstlisting}
\vspace{0.3cm}
\textbf{Output:}\vspace{0.2cm}
\begin{center}
\begin{tabular}{|l|l|}
\hline
\textbf{rollno} & \textbf{name} \\ \hline
92005010 & Ravi Kumar                     \\ \hline
92005011 & Neha Patil                     \\ \hline
\end{tabular}
\end{center}
\newpage
\subsection*{Q3: Students (rollno, name) in ELE dept NOT registered for course EE101}
\begin{lstlisting}[language=SQL]
SELECT
  rollno,
  name
FROM
  students
WHERE
  TRIM(deptcode) = 'ELE'
  AND rollno NOT IN (
    SELECT
      crs_rollno
    FROM
      crs_regd
    WHERE
      TRIM(crs_cd) = 'EE101'
  );
\end{lstlisting}
\vspace{0.3cm}
\textbf{Output:}\vspace{0.2cm}
\begin{center}
\begin{tabular}{|l|l|}
\hline
\textbf{rollno} & \textbf{name} \\ \hline
92005012 & Vikram Singh                   \\ \hline
92005013 & Ananya Roy                     \\ \hline
\end{tabular}
\end{center}
\newpage
\subsection*{Q4: Students registered for BOTH 'DBMS' and 'OS'}
\begin{lstlisting}[language=SQL]
SELECT
  s.name
FROM
  students s
  JOIN crs_regd cr1 ON s.rollno = cr1.crs_rollno
  JOIN crs_offrd co1 ON cr1.crs_cd = co1.crs_code
  JOIN crs_regd cr2 ON s.rollno = cr2.crs_rollno
  JOIN crs_offrd co2 ON cr2.crs_cd = co2.crs_code
WHERE
  TRIM(co1.crs_name) = 'DBMS'
  AND TRIM(co2.crs_name) = 'OS';
\end{lstlisting}
\vspace{0.3cm}
\textbf{Output:}\vspace{0.2cm}
\begin{center}
\begin{tabular}{|l|}
\hline
\textbf{name} \\ \hline
Riddhi Agarwal                 \\ \hline
Priya Verma                    \\ \hline
Sneha Gupta                    \\ \hline
Suresh Nair                    \\ \hline
\end{tabular}
\end{center}
\newpage
\subsection*{Q5: Faculty who offered either 'MIS' or 'Software Engg.'}
\begin{lstlisting}[language=SQL]
SELECT
  f.fac_name
FROM
  faculty f
  JOIN crs_offrd co ON f.fac_code = co.crs_fac_cd
WHERE
  TRIM(co.crs_name) = 'MIS'
  OR TRIM(co.crs_name) = 'Software Engg.';
\end{lstlisting}
\vspace{0.3cm}
\textbf{Output:}\vspace{0.2cm}
\begin{center}
\begin{tabular}{|l|}
\hline
\textbf{fac\_name} \\ \hline
Dr. D.B. Phatak                \\ \hline
Dr. N.L. Sarda                 \\ \hline
\end{tabular}
\end{center}
\newpage
\subsection*{Q6: Faculty who offered 'MIS' but NOT 'Software Engg.'}
\begin{lstlisting}[language=SQL]
SELECT
  f.fac_name
FROM
  faculty f
  JOIN crs_offrd co ON f.fac_code = co.crs_fac_cd
WHERE
  TRIM(co.crs_name) = 'MIS'
  AND f.fac_code NOT IN (
    SELECT
      crs_fac_cd
    FROM
      crs_offrd
    WHERE
      TRIM(crs_name) = 'Software Engg.'
  );
\end{lstlisting}
\vspace{0.3cm}
\textbf{Output:}\vspace{0.2cm}
\begin{center}
\begin{tabular}{|l|}
\hline
\textbf{fac\_name} \\ \hline
Dr. N.L. Sarda                 \\ \hline
\end{tabular}
\end{center}
\newpage
\subsection*{Q7: Students in each hostel not registered for any course}
\begin{lstlisting}[language=SQL]
SELECT
  s.hostel,
  s.rollno,
  s.name
FROM
  students s
WHERE
  s.rollno NOT IN (
    SELECT
      crs_rollno
    FROM
      crs_regd
  );
\end{lstlisting}
\vspace{0.3cm}
\textbf{Output:}\vspace{0.2cm}
\begin{center}
\begin{tabular}{|l|l|l|}
\hline
\textbf{hostel} & \textbf{rollno} & \textbf{name} \\ \hline
4 & 92005013 & Ananya Roy                     \\ \hline
9 & 92005105 & Meera Iyer                     \\ \hline
\end{tabular}
\end{center}
\newpage
\subsection*{Q8: Students in ELE dept OR registered for course CS101}
\begin{lstlisting}[language=SQL]
SELECT
  rollno,
  name
FROM
  students
WHERE
  TRIM(deptcode) = 'ELE'
UNION
SELECT
  s.rollno,
  s.name
FROM
  students s
  JOIN crs_regd cr ON s.rollno = cr.crs_rollno
WHERE
  TRIM(cr.crs_cd) = 'CS101';
\end{lstlisting}
\vspace{0.3cm}
\textbf{Output:}\vspace{0.2cm}
\begin{center}
\begin{tabular}{|l|l|}
\hline
\textbf{rollno} & \textbf{name} \\ \hline
92005001 & Riddhi Agarwal                 \\ \hline
92005002 & Priya Verma                    \\ \hline
92005003 & Rahul Das                      \\ \hline
92005004 & Sneha Gupta                    \\ \hline
92005010 & Ravi Kumar                     \\ \hline
92005011 & Neha Patil                     \\ \hline
92005012 & Vikram Singh                   \\ \hline
92005013 & Ananya Roy                     \\ \hline
92005100 & Suresh Nair                    \\ \hline
92005101 & Kavita Jain                    \\ \hline
92005104 & Arun Bhat                      \\ \hline
\end{tabular}
\end{center}
\newpage
\subsection*{Q9: Students registered for ALL courses}
\begin{lstlisting}[language=SQL]
SELECT
  s.rollno,
  s.name
FROM
  students s
  JOIN crs_regd cr ON s.rollno = cr.crs_rollno
GROUP BY
  s.rollno,
  s.name
HAVING
  COUNT(DISTINCT cr.crs_cd) = (
    SELECT
      COUNT(*)
    FROM
      crs_offrd
  );
\end{lstlisting}
\vspace{0.3cm}
\textbf{Output:}\vspace{0.2cm}
\begin{center}
\begin{tabular}{|l|l|}
\hline
\textbf{rollno} & \textbf{name} \\ \hline
92005001 & Riddhi Agarwal                 \\ \hline
\end{tabular}
\end{center}
\newpage
\subsection*{Q10: Grace Marks +5 in 'DBMS' for students who scored < 50}
\begin{lstlisting}[language=SQL]
UPDATE crs_regd
SET
  marks = marks + 5
WHERE
  marks < 50
  AND crs_cd IN (
    SELECT
      crs_code
    FROM
      crs_offrd
    WHERE
      TRIM(crs_name) = 'DBMS'
  );
\end{lstlisting}
\vspace{0.3cm}
\textit{Statement executed successfully.}\vspace{0.2cm}
\newpage

\end{document}
